\chapter{Runtime and Framework Strategy}
\label{ch:runtime}

This chapter comes before the individual tools because the runtime constraint
shapes every decision that follows. The two interfaces the tool depends on most,
the GENESYS Client API and Roslyn, want incompatible hosts, and the architecture
is largely a response to that.

\section{The Constraint}

The GENESYS Client API assemblies target .NET Framework 4.8. Any code calling
them must run on that framework, and on a machine with a licensed GENESYS
installation present.

Roslyn is available for .NET Framework as well. The core analysis packages,
\code{Microsoft.CodeAnalysis.CSharp} and \code{Microsoft.CodeAnalysis.Common},
ship a \texttt{netstandard2.0} target alongside their modern .NET targets, and
\texttt{netstandard2.0} is loadable from .NET Framework 4.6.1 and above. On
targeting alone, both could sit in the same project.

So the split is not forced by targeting, and it is worth separating two different
reasons for it, because they apply to two different configurations.

\paragraph{For the committed deliverable, the split is a capability decision.}
The tool runs as its own program rather than inside GENESYS, so nothing prevents
Roslyn being referenced from the .NET Framework half. What a modern runtime buys
is the capability that matters most to the analysis:
\code{Microsoft.CodeAnalysis.Workspaces.MSBuild} allows a whole solution to be
opened and resolved rather than individual files, and it behaves far better on a
current .NET runtime than on .NET Framework, where MSBuild assembly resolution is
a known source of difficulty. A current Roslyn also understands recent C\# syntax
while still analysing older code correctly, and carries several years of
performance work.

\paragraph{For the stretch configuration, the split becomes mandatory.}
The requirements list running the tool from the GENESYS interface as a stretch
goal. GENESYS loads its own copy of Roslyn into its process in order to compile
the C\# and Visual Basic scripts it supports, and two versions of the same
assembly cannot coexist in one process. Code running inside GENESYS would
therefore be bound to the version GENESYS ships, which is several years old, or
would break the host's scripting. A fault in the analysis code would also
terminate GENESYS, which the requirements forbid.


\section{The Decision}

The tool is built as two programs in one solution, each on the runtime that suits
it, with a shared contract library between them.

\begin{table}[htbp]
	\centering
	\caption{Projects and their target frameworks}\label{tab:frameworks}
	\begin{tabular}{L{0.26\textwidth}L{0.18\textwidth}L{0.46\textwidth}}
\toprule
	\textbf{Project} & \textbf{Target} & \textbf{Reason} \\
	\midrule
\code{Pentateuch.Genesys} & .NET Framework 4.8 & Must match the Client API
	assemblies. Holds the repository client and the model mapping. \\
	\addlinespace
	\code{Pentateuch.Translator} & .NET 10 & Free to use a current Roslyn, since
	it never shares a process with GENESYS. Holds extraction and generation. \\
	\addlinespace
	\code{Pentateuch.Absyn} & .NET Standard 2.0 & Referenced by both. Holds the
	shared model classes that represent a translation in progress. \\
	\bottomrule
	\end{tabular}
\end{table}

\section{Why .NET Standard 2.0}

.NET Standard is a specification of the API surface a library may use, rather
than a runtime. A library compiled against it can be referenced from any runtime
that implements it, which makes it the only way to share compiled code between
.NET Framework and modern .NET.

Version 2.0 specifically, because it is the highest version .NET Framework
supports. .NET Standard 2.1 was never implemented by .NET Framework and never
will be, so targeting it would exclude the GENESYS-facing half entirely. Version
2.0 is therefore not a preference but the only option that spans both sides.

What this buys the project is a single definition of the abstract syntax.
\code{Pentateuch.Absyn} holds the model classes that represent a translation in
progress, compiled once and referenced by both programs, so the two halves cannot
drift apart in their understanding of the data. Without it, each side would need
its own copy of the same type definitions, kept in step by hand.

The cost is that the shared library is restricted to the .NET Standard 2.0 API
surface, which excludes language and library features added since. In practice
this is a mild restriction for what the library contains, since it holds data
types and no behaviour.

\section{Communication Between the Two Programs}

Because the two programs do not share a process, they cannot pass object
references and must exchange serialised data. \code{Pentateuch.Absyn} defines
what that data is; the mechanism for moving it is a separate decision and has not
yet been made.

At the time of writing the two programs run independently, each exercising its
own interface. \code{Pentateuch.Genesys} connects to a repository and prints a
function's behaviour diagram, and \code{Pentateuch.Translator} generates classes
from a \code{schema.json} file and writes the result to disk. Neither yet calls
the other, so the file-based exchange visible in the translator's test is a test
fixture rather than the chosen transport.

Three candidates remain open: serialised messages over standard input and output,
files written to an agreed location, and named pipes. Standard streams are the
simplest to implement and debug and impose no cleanup burden; files survive the
run and are easy to inspect after the fact; named pipes avoid the filesystem
entirely but add lifetime management. The decision should be recorded here once
made, since it is the one remaining architectural choice in
Table~\ref{tab:frameworks}.

\section{Consequences}

Three things follow from this arrangement, and they are worth stating because
they are frequently mistaken for accidents of implementation.

\paragraph{Most of the tool builds without a GENESYS licence.}
\code{Pentateuch.Translator} and \code{Pentateuch.Absyn} have no dependency on
the vendor assemblies, so they compile and run on any machine with the .NET 10
SDK. Only \code{Pentateuch.Genesys} needs an installation.

The vendor assemblies are referenced from each developer's own installation
rather than committed to the repository, since they are licensed. A
\texttt{GenesysDir} environment variable overrides the default install path for
anyone who installed elsewhere.

\paragraph{The analysis half can use a current Roslyn.}
Solution-level loading, current language support and several years of performance
work are all available, none of which would be if the analysis ran inside
GENESYS.

\paragraph{The split is not reversible later.}
Merging the two halves would reintroduce the assembly conflict, so the boundary
has to exist from the first commit rather than being added when a conflict
appears.
